% Options for packages loaded elsewhere
\PassOptionsToPackage{unicode}{hyperref}
\PassOptionsToPackage{hyphens}{url}
\documentclass[
  11pt,
]{article}
\usepackage{xcolor}
\usepackage[margin=1in]{geometry}
\usepackage{amsmath,amssymb}
\setcounter{secnumdepth}{-\maxdimen} % remove section numbering
\usepackage{iftex}
\ifPDFTeX
  \usepackage[T1]{fontenc}
  \usepackage[utf8]{inputenc}
  \usepackage{textcomp} % provide euro and other symbols
\else % if luatex or xetex
  \usepackage{unicode-math} % this also loads fontspec
  \defaultfontfeatures{Scale=MatchLowercase}
  \defaultfontfeatures[\rmfamily]{Ligatures=TeX,Scale=1}
\fi
\usepackage{lmodern}
\ifPDFTeX\else
  % xetex/luatex font selection
\fi
% Use upquote if available, for straight quotes in verbatim environments
\IfFileExists{upquote.sty}{\usepackage{upquote}}{}
\IfFileExists{microtype.sty}{% use microtype if available
  \usepackage[]{microtype}
  \UseMicrotypeSet[protrusion]{basicmath} % disable protrusion for tt fonts
}{}
\makeatletter
\@ifundefined{KOMAClassName}{% if non-KOMA class
  \IfFileExists{parskip.sty}{%
    \usepackage{parskip}
  }{% else
    \setlength{\parindent}{0pt}
    \setlength{\parskip}{6pt plus 2pt minus 1pt}}
}{% if KOMA class
  \KOMAoptions{parskip=half}}
\makeatother
\usepackage{longtable,booktabs,array}
\newcounter{none} % for unnumbered tables
\usepackage{calc} % for calculating minipage widths
% Correct order of tables after \paragraph or \subparagraph
\usepackage{etoolbox}
\makeatletter
\patchcmd\longtable{\par}{\if@noskipsec\mbox{}\fi\par}{}{}
\makeatother
% Allow footnotes in longtable head/foot
\IfFileExists{footnotehyper.sty}{\usepackage{footnotehyper}}{\usepackage{footnote}}
\makesavenoteenv{longtable}
\setlength{\emergencystretch}{3em} % prevent overfull lines
\providecommand{\tightlist}{%
  \setlength{\itemsep}{0pt}\setlength{\parskip}{0pt}}
\usepackage{amssymb}
\usepackage{microtype}
\setlength{\emergencystretch}{3em}
% Allow long inline-code paths/identifiers to break at separators
% (heading-safe: only redefines behaviour of these chars in \ttfamily,
% no fragile macros in moving arguments).
\usepackage{newunicodechar}
\newunicodechar{§}{\S}
\newunicodechar{²}{\textsuperscript{2}}
\newunicodechar{·}{\ensuremath{\cdot}}
\newunicodechar{ä}{\"{a}}
\newunicodechar{η}{\ensuremath{\eta}}
\newunicodechar{ι}{\ensuremath{\iota}}
\newunicodechar{π}{\ensuremath{\pi}}
\newunicodechar{ᵒ}{\textsuperscript{o}}
\newunicodechar{ᵖ}{\textsuperscript{p}}
\newunicodechar{ᶜ}{\textsuperscript{c}}
\newunicodechar{–}{--}
\newunicodechar{—}{---}
\newunicodechar{₁}{\textsubscript{1}}
\newunicodechar{₂}{\textsubscript{2}}
\newunicodechar{ℤ}{\ensuremath{\mathbb{Z}}}
\newunicodechar{→}{\ensuremath{\rightarrow}}
\newunicodechar{↦}{\ensuremath{\mapsto}}
\newunicodechar{∈}{\ensuremath{\in}}
\newunicodechar{∩}{\ensuremath{\cap}}
\newunicodechar{≅}{\ensuremath{\cong}}
\newunicodechar{≠}{\ensuremath{\neq}}
\newunicodechar{≤}{\ensuremath{\leq}}
\newunicodechar{≥}{\ensuremath{\geq}}
\newunicodechar{⊄}{\ensuremath{\not\subset}}
\newunicodechar{⊓}{\ensuremath{\sqcap}}
\newunicodechar{⊥}{\ensuremath{\bot}}
\newunicodechar{⋙}{\ensuremath{\ggg}}
\newunicodechar{⟨}{\ensuremath{\langle}}
\newunicodechar{⟩}{\ensuremath{\rangle}}
\newunicodechar{⟶}{\ensuremath{\longrightarrow}}
\newunicodechar{⥤}{\ensuremath{\Rightarrow}}
\newunicodechar{𝟭}{\ensuremath{\mathbf{1}}}
\usepackage{bookmark}
\IfFileExists{xurl.sty}{\usepackage{xurl}}{} % add URL line breaks if available
\urlstyle{same}
\hypersetup{
  hidelinks,
  pdfcreator={LaTeX via pandoc}}

\author{}
\date{}

\begin{document}

\section{A Kernel-Checked Music Instance of the Four-Position Partition,
and One Question for Categorical Music
Theory}\label{a-kernel-checked-music-instance-of-the-four-position-partition-and-one-question-for-categorical-music-theory}

\textbf{Author:} Chris Brink (FalseWork) \textbf{Date:} June 2026
\textbf{Status:} Self-contained technical note. The Lean results below
are kernel-checked against Mathlib4 (Lean \texttt{v4.30.0-rc2}); the
single open item (§5) is the intended object of specialist
correspondence. \textbf{Audience:} A categorical music theorist
(Mazzola's denotator/topos framework; Tymoczko's groupoid framework). It
assumes elementary topos theory and pitch-class set theory but nothing
about the surrounding FalseWork program.

\begin{center}\rule{0.5\linewidth}{0.5pt}\end{center}

\subsection{0. One-paragraph summary}\label{one-paragraph-summary}

There is a theorem --- \emph{the four-position partition} --- that says:
in any elementary topos with a non-trivial distinction structure
\texttt{(D,\ η,\ ι)}, the morphisms into a fixed object \texttt{Y} fall
into exactly four cells determined by where \texttt{Im(D\ f)} sits in
the Heyting algebra \texttt{Sub(D\ Y)} relative to the kernel image
\texttt{Im(η\_Y)}. The cells are empty-or-not in a way that is
controlled entirely by Heyting non-regularity. This note reports that
the theorem instantiates \textbf{non-vacuously and kernel-checked} on a
small, entirely standard piece of music mathematics --- the subgroup
lattice of \texttt{ℤ/12}, i.e.~the lattice of transposition-symmetric
pitch-class sets --- with the four cells landing on musically named
objects (tritone, augmented triad, diminished-seventh,
whole-tone/chromatic).

Two features make this more than an illustration. First, the kernel is
\textbf{not a free parameter}: the tritone \texttt{⟨6⟩} is not chosen
and then shown to have four cells around it --- it is the kernel image
\texttt{Im(η)} of a concrete idempotent closure operator
(\texttt{tritoneClosure}, §4), so the four-cell witness is a
\emph{consequence of the operator} rather than a construction built
around a stipulated kernel. Second, the tritone is the \textbf{unique}
kernel in this lattice at which all four cells are non-vacuous: an
exhaustive check over all six possible kernels shows every other choice
degenerates. So this is not ``here is one instantiation'' but ``here is
the only instantiation the lattice admits, and its kernel is forced by
an operator.'' Everything is finite and machine-verified.

The one thing that is \textbf{not} done --- and the reason for writing
to a specialist --- is whether this finite instance is a \emph{slice of
a canonical music topos} (a denotator topos, a presheaf topos on a
Tonnetz groupoid) rather than a bespoke construction. That is §5. If it
is, the four-cell classification is a structural consequence of the
music topos rather than a framework that happens to agree with it.

\begin{center}\rule{0.5\linewidth}{0.5pt}\end{center}

\subsection{1. The theorem being
instantiated}\label{the-theorem-being-instantiated}

Let \texttt{T} be an elementary topos. A \textbf{distinction structure}
on \texttt{T} is an endofunctor \texttt{D\ :\ T\ ⥤\ T}, a unit
\texttt{η\ :\ 𝟭\ ⟶\ D}, and an idempotency iso
\texttt{ι\ :\ D\ ⋙\ D\ ≅\ D} satisfying a coherence law (Spencer-Brown's
``calling'': marking twice equals marking once). It is
\textbf{non-trivial} when \texttt{η} is not a natural iso. For a
morphism \texttt{f\ :\ X\ ⟶\ Y}, write
\texttt{img\ :=\ Im(D\ f)\ ∈\ Sub(D\ Y)} and
\texttt{a\ :=\ Im(η\_Y)\ ∈\ Sub(D\ Y)} (the \emph{kernel image}). The
Heyting algebra \texttt{Sub(D\ Y)} then partitions \texttt{f} (for
\texttt{img\ ≠\ ⊥}) into exactly one of:

{\def\LTcaptype{none} % do not increment counter
\begin{longtable}[]{@{}ll@{}}
\toprule\noalign{}
Cell & Condition on \texttt{img} \\
\midrule\noalign{}
\endhead
\bottomrule\noalign{}
\endlastfoot
Infrastructure & \texttt{img\ ≤\ a} \\
Refusal & \texttt{img\ ≤\ aᶜ} \\
Exploitation & \texttt{img\ ≤\ aᶜᶜ} and \texttt{img\ ⊄\ a} \\
Distribution & \texttt{img\ ⊓\ a\ ≠\ ⊥} and \texttt{img\ ⊓\ aᶜ\ ≠\ ⊥} \\
\end{longtable}
}

The Exploitation cell is empty exactly when \texttt{Sub(D\ Y)} is
Boolean at \texttt{a} (\texttt{aᶜᶜ\ =\ a}); it is the cell that exists
\emph{only} because intuitionistic double negation can be strict.

This is \texttt{four\_position\_partition}, kernel-checked in Lean
(audit: depends only on \texttt{propext}, \texttt{Classical.choice},
\texttt{Quot.sound}; no \texttt{sorry}). The Heyting core is isolated as
\texttt{lattice\_four\_position\_partition}: for any Heyting algebra
\texttt{H}, any \texttt{a\ :\ H}, and any \texttt{x\ ≠\ ⊥}, exactly one
of the four conditions holds. The full statement and proof are in the
partition paper (\texttt{paper.md}).

The bridge from ``an idempotent operation'' to ``a distinction
structure'' is also kernel-checked:
\texttt{DistinctionStructure.ofIdempotentMonad} (Remark 5.5 of the
paper) turns any idempotent monad --- equivalently any reflective
subcategory, equivalently (at the subobject level) any closure operator
on \texttt{Sub(D\ Y)} --- into a distinction structure. This is the
hinge the music instance hangs on.

\subsection{\texorpdfstring{2. The substrate: the subgroup lattice of
\texttt{ℤ/12}}{2. The substrate: the subgroup lattice of ℤ/12}}\label{the-substrate-the-subgroup-lattice-of-ux212412}

Take the six subgroups of \texttt{ℤ/12}. Ordered by inclusion they form
the divisor lattice of 12:

\begin{verbatim}
              ℤ/12               (full chromatic)
             /     \
          ⟨2⟩       ⟨3⟩
        (whole-tone) (dim 7th)
          |   \     /
          |    \   /
         ⟨6⟩    ⟨4⟩
       (tritone) (aug triad)
             \   /
             ⟨0⟩                  (trivial)
\end{verbatim}

Each subgroup \emph{is} a transposition-symmetric pitch-class set:
\texttt{⟨6⟩\ =\ \{0,6\}} (tritone), \texttt{⟨4⟩\ =\ \{0,4,8\}}
(augmented triad), \texttt{⟨3⟩\ =\ \{0,3,6,9\}} (diminished seventh),
\texttt{⟨2⟩\ =\ \{0,2,4,6,8,10\}} (whole-tone hexachord). These are
Messiaen's modes of limited transposition / the symmetric sets of
pitch-class set theory --- they predate and are independent of any
specific categorical framework (Tymoczko's, Mazzola's, Lewin's).

Two facts make this the right toy substrate:

\begin{enumerate}
\def\labelenumi{\arabic{enumi}.}
\tightlist
\item
  \textbf{It is a non-Boolean Heyting algebra.} Divisor lattices are
  distributive (hence Heyting); \texttt{12\ =\ 2²·3} is not squarefree,
  so the lattice is not Boolean. Concretely the tritone is
  \textbf{non-regular}: \texttt{⟨6⟩ᶜᶜ\ =\ ⟨3⟩\ ≠\ ⟨6⟩}. Kernel-checked:
  \texttt{Div12.tritone\_non\_regular}.
\item
  \textbf{The literal pitch-class realization is now machine-checked,
  not asserted.} The map \texttt{pcset\ :\ Div12\ →\ Finset\ (ℤ/12)}
  sending each lattice element to its actual pitch-class set is proved
  injective, valued in additive subgroups (closed under \texttt{0},
  \texttt{+}, negation), order-reflecting (lattice order = subgroup
  inclusion), and meet-preserving
  (\texttt{pcset\ (a\ ⊓\ b)\ =\ pcset\ a\ ∩\ pcset\ b}). Bundled as
  \texttt{Div12.pcset\_realizes\_subgroup\_lattice}.
\end{enumerate}

\subsection{3. The four cells, instantiated and
named}\label{the-four-cells-instantiated-and-named}

\textbf{The tritone is the unique kernel that works.} Before fixing it:
an exhaustive check over all six possible kernels in this lattice
(computational companion,
\texttt{wolfram/music-anchor/layer-t-d-checks.wl} §5) shows the tritone
\texttt{⟨6⟩} is the \emph{only} kernel at which all four cells are
non-vacuously inhabited. Every other choice degenerates --- the other
non-regular element (whole-tone \texttt{⟨2⟩}) has \texttt{⟨2⟩ᶜ\ =\ ⊥},
collapsing Refusal and Distribution; the four regular elements collapse
further. So the kernel below is not one choice among six but the forced
one, a structural fact about the lattice rather than a tuning of the
example.

Fix the kernel \texttt{a\ =\ ⟨6⟩} (tritone). Then \texttt{aᶜ\ =\ ⟨4⟩}
(augmented triad) and \texttt{aᶜᶜ\ =\ ⟨3⟩} (diminished seventh). The
partition reads off the lattice as:

{\def\LTcaptype{none} % do not increment counter
\begin{longtable}[]{@{}
  >{\raggedright\arraybackslash}p{(\linewidth - 4\tabcolsep) * \real{0.1500}}
  >{\raggedright\arraybackslash}p{(\linewidth - 4\tabcolsep) * \real{0.3500}}
  >{\raggedright\arraybackslash}p{(\linewidth - 4\tabcolsep) * \real{0.5000}}@{}}
\toprule\noalign{}
\begin{minipage}[b]{\linewidth}\raggedright
Cell
\end{minipage} & \begin{minipage}[b]{\linewidth}\raggedright
Inhabited by
\end{minipage} & \begin{minipage}[b]{\linewidth}\raggedright
Pitch-class object
\end{minipage} \\
\midrule\noalign{}
\endhead
\bottomrule\noalign{}
\endlastfoot
Infrastructure (\texttt{x\ ≤\ a}) & \texttt{⟨6⟩} & tritone (the kernel
itself) \\
Refusal (\texttt{x\ ≤\ aᶜ}) & \texttt{⟨4⟩} & augmented triad \\
Exploitation (\texttt{x\ ≤\ aᶜᶜ}, \texttt{x\ ⊄\ a}) & \texttt{⟨3⟩} &
diminished seventh (the closure-residue) \\
Distribution (straddle) & \texttt{⟨2⟩}, \texttt{ℤ/12} & whole-tone
hexachord, full chromatic \\
\end{longtable}
}

All four cells inhabited, kernel-checked as
\texttt{Div12.music\_anchor\_witness} (uniqueness of this kernel stated
at the top of the section).

\textbf{A correction worth stating up front}, because it is the obvious
first guess and it is wrong: one is tempted to map Tymoczko's three
harmonic fields --- diatonic, symmetric, chromatic --- onto
Infrastructure/Exploitation/Refusal via a \emph{diatonic-closure}
distinction operator (close a pitch-class set to the smallest diatonic
scale containing it). We tried exactly this (memo §11). The
diatonic-closure Moore lattice on \texttt{ℤ/12} (92 closed sets) is
\textbf{not distributive}, hence not Heyting, with two explicit
counterexamples to the Heyting identity. So the partition theorem does
not apply to that construction at all. The working instance is the
subgroup lattice with the \textbf{tritone} as kernel, above --- a
different and cleaner object. Any bridge that routes through diatonic
closure inherits the §11 obstruction; the subgroup route avoids it.

\subsection{4. The distinction operator (this note's new
content)}\label{the-distinction-operator-this-notes-new-content}

The witness in §3 takes \texttt{a\ =\ ⟨6⟩} as \emph{given}. The new
result here is that this kernel is \textbf{not a free parameter}: it is
the kernel image of a concrete idempotent operation.

Define the closure operator \texttt{tritoneClosure\ :\ Div12\ →\ Div12}
whose closed (fixed) points are the Moore family \texttt{\{⟨6⟩,\ ℤ/12\}}
--- i.e.~the closure of \texttt{x} is the least closed element
\texttt{≥\ x}:

\begin{verbatim}
⊥ ↦ ⟨6⟩,  ⟨6⟩ ↦ ⟨6⟩,  ⟨4⟩ ↦ ℤ/12,  ⟨3⟩ ↦ ℤ/12,  ⟨2⟩ ↦ ℤ/12,  ℤ/12 ↦ ℤ/12.
\end{verbatim}

Kernel-checked facts
(\texttt{Div12.tritoneClosure\_is\_distinction\_slice}):

\begin{itemize}
\tightlist
\item
  \textbf{extensive} \texttt{x\ ≤\ tritoneClosure\ x},
  \textbf{monotone}, \textbf{idempotent} --- a bona fide closure
  operator;
\item
  its fixed-point set is exactly \texttt{\{⟨6⟩,\ ℤ/12\}}
  (\texttt{tritoneClosure\_fixedPoints}), meet-closed (a Moore family);
\item
  \textbf{\texttt{tritoneClosure\ ⊥\ =\ ⟨6⟩}}
  (\texttt{tritoneClosure\_bot}): the closure of the bottom element is
  the tritone;
\item
  \texttt{(tritoneClosure\ ⊥)ᶜᶜ\ ≠\ tritoneClosure\ ⊥}
  (\texttt{tritoneClosure\_bot\_non\_regular}): the kernel image is
  non-regular, which is \emph{precisely} the condition that the
  Exploitation cell is non-empty.
\end{itemize}

In the topos shadow this is the picture: a closure operator on
\texttt{Sub(T)(1)} is what an idempotent monad's unit cuts out, and
\texttt{Im(η)} corresponds to the closure of the initial object
\texttt{⊥}. So \texttt{tritoneClosure\ ⊥\ =\ ⟨6⟩} is the lattice-level
statement that \textbf{\texttt{Im(η)} lands on the tritone}, making
\texttt{a\ =\ ⟨6⟩} the kernel image of a genuine distinction operation
rather than a stipulation. This is candidate \texttt{\{2,12\}} --- the
minimal tritone-closing Moore family --- from the enumerated Layer-D
candidate space (memo §12.6; there are exactly four tritone-closing
families).

This closes the lattice-level (Layer-L + Layer-D-slice) story end to
end: a non-Boolean music-derived Heyting algebra, a concrete idempotent
distinction operator on it whose kernel image is non-regular, and the
four cells inhabited at that kernel image --- all kernel-checked.

\subsection{5. The one open question (where a specialist is
needed)}\label{the-one-open-question-where-a-specialist-is-needed}

What is \textbf{not} done is the topos-level lift, and it is the only
thing standing between this and a complete
\texttt{four\_position\_partition} instance on a \emph{music topos}:

\begin{quote}
Is there a canonical elementary topos \texttt{T} from categorical music
theory --- a denotator topos in your framework, or a presheaf topos on a
Tonnetz/pitch-class groupoid --- such that \texttt{Sub\_T(1)} (or
\texttt{Sub\_T(Y)} for a natural \texttt{Y}) is the subgroup lattice of
\texttt{ℤ/12}, and such that \texttt{tritoneClosure} is the
subobject-level trace of an idempotent monad on \texttt{T} (a
sheafification, a reflection onto a sub-topos, a symmetrization)?
\end{quote}

Two existence answers are already in hand from general topos theory, but
both are \emph{generic}, with no music content beyond what the lattice
carries: (T1) \texttt{Sh(L)} for the lattice \texttt{L} as a locale;
(T2) \texttt{Set\^{}\{Pᵒᵖ\}} for \texttt{P} the 3-element poset of
join-irreducibles \texttt{\{⟨6⟩\ \textless{}\ ⟨3⟩,\ ⟨4⟩\}} (Birkhoff),
whose lattice slice is computationally verified. Neither tells us
whether the \emph{music} topos you would actually build realizes this
lattice and this closure operator.

That is the specialist question. Three concrete sub-questions:

\begin{enumerate}
\def\labelenumi{\arabic{enumi}.}
\tightlist
\item
  \textbf{Realization.} Does the denotator/local-composition topos (or a
  presheaf topos on the dihedral pitch-class groupoid \texttt{D₁₂})
  contain an object \texttt{Y} with \texttt{Sub\_T(Y)\ ≅} the subgroup
  lattice of \texttt{ℤ/12}? The subobject lattice of a representable
  presheaf on a groupoid is the lattice of sub-\texttt{Aut}-sets; for
  \texttt{ℤ/12} acting on itself this is close, but the precise object
  needs your eye.
\item
  \textbf{The operator.} Is \texttt{tritoneClosure} the subobject-level
  trace of a natural idempotent monad on \texttt{T} --- e.g.~a
  symmetrization (orbit under a chosen symmetry subgroup), or a
  sheafification for a Lawvere--Tierney topology corresponding to
  ``tritone coverage''? Idempotency you can read off; whether it is the
  \emph{intended} distinction operation is a question about your
  framework's semantics, not just its mathematics.
\item
  \textbf{Better \texttt{D}.} If \texttt{tritoneClosure} is the wrong
  operator, is there a better one --- your global-composition /
  denotator-limit construction, say --- whose kernel image still lands
  non-regular, giving a non-empty Exploitation cell with a different
  (perhaps richer) musical reading?
\end{enumerate}

If the answer to (1)+(2) is yes, then the four-cell classification is
not an empirical framework that \emph{agrees} with the categorical
apparatus --- it is a structural consequence of it, obtained through the
partition theorem. If (2) needs (3), that is itself informative about
which distinction operation is canonical for music.

\subsection{6. Scope honesty (what this note does not
claim)}\label{scope-honesty-what-this-note-does-not-claim}

\begin{itemize}
\tightlist
\item
  It does \textbf{not} claim a topos-level instance. §3--§4 are
  lattice-level (Layer L and the Layer-D \emph{slice}); the topos lift
  of §5 is open.
\item
  It does \textbf{not} claim the diatonic/symmetric/chromatic
  tripartition maps onto the cells. The naive (diatonic-closure) form of
  that map fails at the Heyting-distributivity step (§3, memo §11); any
  correct version must be reconstructed on the subgroup substrate and is
  not established here.
\item
  It does \textbf{not} claim a formal identity between the tritone's
  Heyting non-regularity and the Pythagorean comma's irrationality
  \texttt{(3/2)\^{}p\ ≠\ 2\^{}q}. The two are recorded as a
  \textbf{structural identification} (memo §13.1) --- same load-bearing
  role, ``comma as the obstruction to closure'' --- not a constructed
  isomorphism. A locale/\texttt{π₁} bridge connecting the
  lattice-of-subobjects picture to a topology-of-arrows picture is
  plausible but unbuilt.
\item
  The pitch-class objects used are elementary group theory of
  \texttt{ℤ/12} (Messiaen, Forte, Rahn), not a contribution of any
  recent categorical author. The instrument is borrowed; only the
  partition reading of its cells is the framework's.
\end{itemize}

\subsection{7. Reproducing the kernel-checked
claims}\label{reproducing-the-kernel-checked-claims}

All results live in the public repository (\texttt{thefalsework/papers},
\texttt{lean/} directory), Lean \texttt{v4.30.0-rc2} against Mathlib4,
in namespaces \texttt{FalseWork.Positions} and
\texttt{FalseWork.Lattice} (the witnesses in
\texttt{FalseWork.Lattice.Examples.Div12}). By file:

\begin{verbatim}
Positions/Partition.lean
    four_position_partition           (the topos-level theorem)
Positions/SpencerBrown.lean
    DistinctionStructure.ofIdempotentMonad   (Remark 5.5 bridge)
Lattice/FourPositionLattice.lean
    lattice_four_position_partition          (the Heyting core)
Examples/DivisorLattice12.lean            (Layer-L witness)
    tritone_non_regular
    music_anchor_witness
Examples/DivisorLattice12Distinction.lean (this note's results)
    tritoneClosure_is_distinction_slice
    pcset_realizes_subgroup_lattice
    pcset_tritoneClosure_bot
\end{verbatim}

Axiom audit (\texttt{Examples/HeytingTypeInstance.lean}, via
\texttt{\#print\ axioms}): every theorem above depends only on
\texttt{propext}, \texttt{Classical.choice}, \texttt{Quot.sound} ---
several on strictly fewer --- and none on \texttt{sorryAx}.

\begin{center}\rule{0.5\linewidth}{0.5pt}\end{center}

\subsubsection{References}\label{references}

\begin{itemize}
\tightlist
\item
  Brink, C. (2026). \emph{A Four-Position Partition of Morphisms in
  Elementary Topoi with Distinction Structure.} Preprint,
  \texttt{paper.md} (this repo).
\item
  Mazzola, G. (2002). \emph{The Topos of Music.} Birkhäuser.
\item
  Messiaen, O. (1944). \emph{Technique de mon langage musical.}
\item
  Tymoczko, D. (2011). \emph{A Geometry of Music.} OUP. --- (2026). The
  concept of musical space. \emph{Journal of Music Theory} 70(1).
\item
  Music-anchor feasibility memo: \texttt{music-anchor/feasibility.md}
  (§11 diatonic-closure negative result; §12 subgroup route and layered
  status; §13 Tymoczko correspondence).
\end{itemize}

\end{document}
